\documentclass[10pt,a4paper]{amsart}
\usepackage[margin=19mm]{geometry}
\usepackage{amsmath,amssymb,mathtools}
\usepackage[scr=rsfs,cal=euler]{mathalfa}
\usepackage{xfrac}
\usepackage[unicode,hidelinks,bookmarks=false]{hyperref}
\newcommand{\ie}{i.e.\ }
\newcommand{\cf}{cf.\ }
\newcommand{\resp}{respectively\ }
\newcommand{\Subsection}{\subsection}
\providecommand{\nobreakdash}{\nobreak\hbox{-}}
\setlength{\emergencystretch}{2em}
\multlinegap0pt
\usepackage{bm}
\usepackage{bbm}
\usepackage{dsfont}
\usepackage{xspace}
\usepackage{cancel}
\usepackage{mathtools}
\usepackage{amsthm}
\makeatletter
\@ifundefined{theorem}{\newtheorem{theorem}{Theorem}}{}
\makeatother
\title[Model theory of term algebras revisited]{Model theory of term algebras revisited}
\author{Davide Carolillo, Yifan Jia, Bakh Khoussainov, Rizos Sklinos}
\date{Source: arXiv:2602.01819v1 (CC BY 4.0)}
\begin{document}
\maketitle

\noindent\emph{Source and licence.} Model theory of term algebras revisited. Version arXiv:2602.01819v1 (CC BY 4.0), distributed under the Creative Commons Attribution 4.0 International licence, as recorded in the arXiv licence metadata for this exact version and shown on the abstract page https://arxiv.org/abs/2602.01819v1. Full rights notice: licenses/arxiv-2602.01819-CC-BY-4.0.txt.\par\smallskip

\noindent\textit{Editorial scope.} Counted unit: the Theorem whose source label is \texttt{IndecomposableEquivalentElementary} in the article (source0.tex lines 283--285, source label \texttt{IndecomposableEquivalentElementary}), delivered together with the article's own complete original proof of it (source0.tex lines 286--310, one \texttt{proof} environment of 25 lines). No mathematics is written, repaired or re-derived: statement and proof are the article's own text line for line. Required context retained uncounted: none. Every definition, display and label of those lines is retained unchanged. Omitted from this package and not redistributed: 1--282, 311--812. Nothing inside a retained excerpt is omitted or altered: every retained body is delivered exactly as the article's lines are written, so no omission receipt of any kind accompanies this package. Citations: the retained excerpts carry no citation command at all, so no bibliography entry of the article is transcribed and no reference is redistributed. Reference adaptation: none was needed and none was made; every reference inside the delivered unit resolves inside it, so no unresolved reference or citation is produced. Printed slips kept literal: no printed word, symbol, hypothesis or number of the counted statement or of its proof was altered. Layout and class: the article's own class is \texttt{article} with packages including xcolor, leftindex, mathtools, amsmath, qtree, forest, todonotes; the wrapper is the lane's pinned \texttt{amsart} editorial wrapper at 10pt on A4, which restates the theorem-like environments the retained text needs and adds the article's own macro definitions that the retained text needs (none), with extra wrapper packages (bm, bbm, dsfont, xspace, cancel, mathtools). The delivered numbering is the unit's own. Assets: no retained span contains a figure environment, a graphics-inclusion command, a PDF-page inclusion, a table or a TikZ picture; no image file is copied into this package, the package declares no asset dependency and no image file is redistributed with it. Licence: version arXiv:2602.01819v1 of this article is distributed under the Creative Commons Attribution 4.0 International licence, as recorded in the arXiv licence metadata for this exact version and shown on the abstract page https://arxiv.org/abs/2602.01819v1. A full rights notice is delivered at licenses/arxiv-2602.01819-CC-BY-4.0.txt. Characters: every retained excerpt is pure ASCII, so the delivered engine, XeLaTeX with its default Latin Modern text font, reports no missing character.
\begin{theorem}\label{IndecomposableEquivalentElementary}
Let $0\leq m<\omega$. Then the following list of axioms $(A_i)_{i\leq k}$, $(B_{ij})_{0\leq i<j\leq k}$, $(C_t)_{t\neq x}$ and $D_m$ is complete. 
\end{theorem}

\begin{proof}
We will show that if two $\Sigma$-structures, $\mathcal{A}$ and $\mathcal{B}$, satisfy the axioms of the hypothesis, they are elementarily equivalent. Equivalently, we need to show that for each $r<\omega$, Player II has a winning strategy in the $r$-round game $EF_r[\mathcal{A}, \mathcal{B}]$. We recall that unnested atomic formulas in $\Sigma$ are formulas of the form, $x=y, f_0(x_1, \ldots, x_{n_0})=y, \ldots, f_k(x_1,\ldots, x_{n_k})=y$. The strategy is as follows: first of all, by Axiom $D_m$, the structures $\mathcal{A}$ and $\mathcal{B}$ have the same number of indecomposable elements, namely $m$. We choose a "random" correspondence between the two sets of indecomposable elements, unless $m=0$, in which case this first step is not necessary. 

At every round of the game Player II copies the skeleton (of depth that depends on the round) of the choice of Player I (independent of which structure player I chooses) and generates an element in the respective structure by choosing the generating elements according to the isomorphism of the indecomposable elements but also the "inventory" (to be defined) of the previous round.

We define the strategy recursively:
\begin{itemize}
\item At round 1, suppose Player I plays element $a_1$ from $\mathcal{A}$. Player II copies $sk_r(a_1)$ and checks $dc_r(a_1)$ (if $sk_r(a_1)=\bot$, then it copies $sk_l(a)$ for the largest $l\leq r$ for which $sk_l(a_1)\neq\bot$). For any element in $dc_r(a_1)$ which is indecomposable Player II will replace it with the corresponding indecomposable element in $\mathcal{B}$. On the other hand,  for any decomposable element in $dc_r(a_1)$, Player II will choose at random a decomposable element in $\mathcal{B}$. 
At this point, we extend the isomorphism between the indecomposable elements to include the corresponding choices of this round among all elements between $dc_{\leq r}(a_1)$ and $dc_{\leq r}(b_1)$, where $b_1$ is the actual answer of Player II. The inventory of the first round $I_1$ is the data that consists of: $dc_{\leq r}(a_1), dc_{\leq r}(b_1)$ the isomorphism between the indecomposable elements and the isomorphism between $dc_{\leq r}(a_1)$ and $ dc_{\leq r}(b_1)$. Note that the strategy would be the same if Player I had chosen an element from $\mathcal{B}$. Also, it is important that we keep in the inventory all $dc_{\leq r}(a_1)$ and $ dc_{\leq r}(b_1)$ and not just $dc_{r}(a_1)$ and $ dc_{r}(b_1)$.

\item At round $i+1$, suppose Player I plays element $a_{i+1}$ from $\mathcal{A}$. Player II copies the skeleton $sk_{r-(i+1)}(a_i)$ (or the largest skeleton less than $r-(i+1)$ which is not $\bot$) and checks $dc_{r-(i+1)}(a_i)$, he then uses the inventory of the previous move, $I_i$, and chooses generating elements accordingly. The new inventory, $I_{i+1}$, is $I_i$ together with the new elements in $dc_{\leq r-(i+1)}(a_i)$, $dc_{\leq r-(i+1)}(b_i)$, where $b_i$ is the actual answer of Player II for this round, and the isomorphism between them.
\end{itemize}

We now prove that this strategy works, i.e. it is indeed a winning strategy for Player II.

For a warm up we start with the unnested formula $x=y$. Suppose $a_i=a_j$, and without loss of generality we may assume that $i>j$. Then, since the inventory of the $i-1$-round contains $a_j$ and the comprehensive downward closure of an element in later rounds is contained in the comprehensive downward closure of the element in earlier rounds (since the depth decreases as we move forward in the game), we get that $b_i$ admits a construction as $b_j$ (independent of whether Player I played $b_i$ or $a_i$). Hence $b_i=b_j$.   

We move to the case $y=f(x_1, \ldots, x_n)$. Suppose $a_j=f(a_{i_1}, a_{i_2}, \ldots, a_{i_n})$, for some $j\leq r$ and $f$ some function symbol of arity $n$ from $\Sigma$. Notice that $a_j$ cannot appear in $f$. Also for any $a_i$, with $i<j$, $a_i$ belongs to the inventory of the $j-1$-th round and since $sk_1(a_j)=f(x_1,\ldots,x_r)$, we have that $sk_1(b_j)=f(x_1, \ldots,x_r)$. In addition, as before, the comprehensive downward closure of an element in later rounds is contained in the comprehensive downward closure of the element in earlier rounds (since the depth decreases as we move forward in the game). This implies that if $a_1$, for example, appears in $f$, say in the $l$-th argument, then $b_1$ must appear in the $l$-th argument of $f$ in the construction of $b_j$ as an image of $f$ (independent of whether Player I played $a_j$ or $b_j$). 

We, thus, need to show that if $a_i$, for $i>j$, appears in $f$, then $b_i$ must also appear in the same argument of $f$ in the corresponding construction of $b_j$. Suppose $a_i$ is the $l$-th argument of $f$ (it may appear in more than one places but this will not make any difference). The $j-1$ comprehensive downward closure of $a_i$ will be contained in the $I_j$ inventory and it will have a corresponding element $b'$ from the structure $\mathcal{B}$. We now take cases according to whether Player I plays $a_i$ or $b_i$. Assume first that in the $i$-round Player I plays element $a_i$ from $\mathcal{A}$, then $a_i$ is determined (since it appears in the inventory of a previous round) i.e. its construction does not require any free choices. In particular, $b_i$ will be constructed using only elements from the inventory and thus, it must be equal to $b'$ which is the $l$-th argument of $f$ as we wanted. Next assume that Player I plays element $b_i$ from $\mathcal{B}$ and assume $b_i$ is not equal to $b'$. On the other hand, the $(r-i)$-skeleton of $b_i$ and $a_i$ are the same and the $(r-i)$-downward closure of $a_i$ consists of elements in correspondence with the $(r-i)$-downward closure of $b_i$. We claim that  the $(r-i)$-skeleton of $b_i$ must be the same as the $(r-i)$-skeleton of $b'$. Indeed, both must be equal to the $(r-i)$-skeleton of $a_i$ which is equal to the $(r-i)$-skeleton of the $l$-th argument of $f$ in the analysis of $a_j$. Hence, there must be an element in the $(r-i)$-downward closure of $b_i$ that differs from that of $b'$. In particular, since there is a correspondence between the elements of the $(r-i)$-downward closure of $b'$ and those of $a_i$ (given in the $I_j$ inventory), we must have that there is an element in the construction of $a_i$ in the $i$-th move that differs from its construction as the $l$-th argument of $f$ in the analysis of $a_j$ in the $j$-th round, in particular $a_j\neq f(a_{i_1}, a_{i_2}, \ldots, a_{i_n})$, a contradiction.    
    
Notice that it is crucial in the strategy that Player II copies the skeleton of lower and lower depth as the game moves forward. Indeed, if independent of the round Player II kept copying the $r$-skeleton of the choice of Player I, then it can be easily seen that Player I may win.  

%Whenever, at any round of the game, Player I chooses an element from $\mathcal{F}_m$, Player II replies by a copy of that element in $\mathcal{A}$ using the fixed correspondence of indecomposable elements. Likewise, when Player I plays an element of finite height from $\mathcal{A}$, i.e. an element in the subalgebra generated by the indecomposable elements, then Player II copies the element in $\mathcal{F}_m$ using the fixed correspondence. In particular, if Player I plays an indecomposable element (in either structures), then Player II plays the corresponding indecomposable element. If, at round $i$ of the game, Player I plays an element from $\mathcal{A}$ which does not have finite height then Player I starts unfolding it up to $k+1$ steps, for the elements that are still decomposable after $k+1$ steps, it replaces them with elements of height $>k+1$ from $\mathcal{F}_m$, bookkeeping the equality of terms. At round $i+1$,  We show that this is a winning strategy for Player II. 
\end{proof}

\end{document}
